\section{完全可积系统}\label{sec:completely-integrable-systems}

现在集中研究 Hamilton 函数 $H\in C^\infty(T^*M)$ 的 Hamilton 流. 已知 $H$ 沿其 Hamilton 流保持恒定, 因此只需在 $H$ 的水平集上研究此流. 如果还存在其他非平凡函数 $F\in C^\infty(T^*M)$, 使得
\[
  X_H\cdot F=0\iff\{H,F\}=0,
\]
那么还可以进一步简化这一研究.

\begin{RiemannDefinition}\label{def:first-integral-of-hamiltonian}
若 $\{H,F\}=0$, 则称函数 $F\in C^\infty(T^*M)$ 为 $H$ 的\textbf{第一积分}.
\end{RiemannDefinition}

一般而言, 并无理由期待存在 $H$ 本身之外的非平凡第一积分. 在它们确实存在的特殊情形中, 往往还满足一些附加条件.

\begin{RiemannDefinition}\label{def:involution-and-independent-integrals}
对于函数 $F_1,\ldots,F_m\in C^\infty(T^*M)$:
\begin{enumerate}[label=(\roman*),leftmargin=*]
\item 若 $\{F_i,F_j\}=0$ ($i,j=1,\ldots,m$), 则称它们\textbf{对合};
\item 若 $(\mathrm{d}F_1)_\alpha,\ldots,(\mathrm{d}F_m)_\alpha\in T_\alpha^*(T^*M)$ 是线性无关的余向量, 则称它们在 $\alpha\in T^*M$ 处\textbf{独立}.
\end{enumerate}
\end{RiemannDefinition}

\begin{RiemannProposition}\label{prop:dimension-bound-for-independent-integrals-in-involution}
若 $F_1,\ldots,F_m\in C^\infty(T^*M)$ 对合且在某点 $\alpha\in T^*M$ 处独立, 则 $m\leq n$.
\end{RiemannProposition}

\begin{RiemannProof}
Exercise~\ref{ex:completely-integrable-systems}(\ref{ex:prove-involutive-integral-dimension-bound}).
\end{RiemannProof}

达到上界 $m=n$ 的情形尤其值得关注.

\begin{RiemannDefinition}\label{def:completely-integrable-hamiltonian}
若存在 $n$ 个对合的第一积分 $F_1,\ldots,F_n$, 且它们在稠密开集 $U\subset T^*M$ 上独立, 则称 Hamilton 函数 $H$ \textbf{完全可积}.
\end{RiemannDefinition}

\begin{RiemannExample}\label{ex:complete-integrability-in-dimension-one-and-central-fields}
\begin{enumerate}[label=(\arabic*),ref=\arabic*,leftmargin=*]
\item 若 $M$ 是 1 维的, 且在 $T^*M$ 的一个稠密开集上有 $\mathrm{d}H\ne0$, 则 $H$ 完全可积.
\item (\emph{中心力场中的粒子}) 回顾 Example~\ref{ex:central-fields-and-spherical-christoffel-symbols} 中质量为 $m>0$ 的粒子在中心力场中运动的情形. 相应的 Lagrange 函数为
\[
  L(r,\theta,v^r,v^\theta)=\frac12 m\bigl[(v^r)^2+r^2(v^\theta)^2\bigr]-u(r),
\]
所以 Legendre 变换由
\[
  p_r=\frac{\partial L}{\partial v^r}=mv^r
  \qquad\text{和}\qquad
  p_r=\frac{\partial L}{\partial v^\theta}=mr^2v^\theta
\]
给出. 于是 Hamilton 函数为
\[
  H(r,\theta,p_r,p_\theta)=\frac{p_r^2}{2m}+\frac{p_\theta^2}{2mr^2}+u(r).
\]
由 Hamilton 方程,
\[
  \dot p_\theta=-\frac{\partial H}{\partial\theta}=0,
\]
所以 $p_\theta$ 是一个第一积分. 由于
\[
  \mathrm{d}H=\left(-\frac{p_\theta^2}{mr^3}+u'(r)\right)\mathrm{d}r
  +\frac{p_r}{m}\,\mathrm{d}p_r+\frac{p_\theta}{mr^2}\,\mathrm{d}p_\theta,
\]
可见 $\mathrm{d}H$ 与 $\mathrm{d}p_\theta$ 在 $T^*\mathbb{R}^2$ 的一个稠密开集上独立. 这个开集由极坐标 $(r,\theta,p_r,p_\theta)$ 有定义且不满足
\[
  u'(r)-\frac{p_\theta^2}{mr^3}=p_r=0
\]
的点构成 (即不在圆轨道上的点 --- 参见 Exercise~\ref{ex:completely-integrable-systems}(\ref{ex:stability-of-central-field-circular-orbits})). 因此这个 Hamilton 函数完全可积.
\end{enumerate}
\end{RiemannExample}

\begin{RiemannProposition}\label{prop:invariant-level-sets-for-integrable-system}
设 $H$ 为完全可积 Hamilton 函数, 具有在稠密开集 $U\subset T^*M$ 上独立的对合第一积分 $F_1,\ldots,F_n$, 且 $X_{F_1},\ldots,X_{F_n}$ 在 $U$ 上完备. 则每个非空水平集
\[
  L_f:=\{\alpha\in U\mid F_1(\alpha)=f_1,\ldots,F_n(\alpha)=f_n\}
\]
都是 $n$ 维子流形, 在 $H$ 的 Hamilton 流下不变, 并容许 $\mathbb{R}^n$ 的局部自由作用, 该作用在每个连通分支上传递.
\end{RiemannProposition}

\begin{RiemannProof}
$U$ 中所有点都是映射 $F:U\to\mathbb{R}^n$ 的正则点, 其中 $F(\alpha)=(F_1(\alpha),\ldots,F_n(\alpha))$; 因此所有非空水平集 $L_f:=F^{-1}(f)$ 都是 $n$ 维子流形.

由于对 $i=1,\ldots,n$ 有 $X_H\cdot F_i=0$, 水平集 $L_f$ 在 $X_H$ 的流下不变. 此外, 有 $X_{F_i}\cdot F_j=\{F_i,F_j\}=0$, 因而这些水平集在 $X_{F_i}$ 的流下也不变. 又由于 $[X_{F_i},X_{F_j}]=X_{\{F_i,F_j\}}=0$, 这些流彼此交换 (参见第~\ref{chap:differentiable-manifolds}~章的 Theorem~\ref{thm:commuting-flows-and-lie-brackets}).

考虑由下式给出的映射 $A:\mathbb{R}^n\times L_f\to L_f$:
\[
  A(t_1,\ldots,t_n,\alpha)=(\psi_{1,t_1}\circ\cdots\circ\psi_{n,t_n})(\alpha),
\]
其中 $\psi_{i,t}:L_f\to L_f$ 是 $X_{F_i}$ 的流. 由于这些流彼此交换, 该映射定义了 $\mathbb{R}^n$ 在 $L_f$ 上的作用. 另一方面, 对每个 $\alpha\in L_f$, 由 $A_\alpha(t_1,\ldots,t_n)=A(t_1,\ldots,t_n,\alpha)$ 给出的映射 $A_\alpha:\mathbb{R}^n\to L_f$ 在原点处是局部微分同胚, 因为
\[
  (\mathrm{d}A_\alpha)_0(e_i)=\left.\frac{\mathrm{d}}{\mathrm{d}t}\right|_{t=0}
    \psi_{i,t}(\alpha)=(X_{F_i})_\alpha
\]
且向量场 $X_{F_i}$ 线性无关. 因此该作用是局部自由的 (即对每个点 $\alpha\in L_f$, 都存在 $0$ 的一个开邻域 $U\subset\mathbb{R}^n$, 使对所有 $t\in U\setminus\{0\}$ 都有 $A(t,\alpha)\ne\alpha$), 也就是说, 稳定子群是离散的. 此外, 该作用是局部传递的 (即每个点 $\alpha\in L_f$ 都有开邻域 $V\subset L_f$, 使每个 $\beta\in V$ 都形如 $\beta=A(t,\alpha)$, 其中 $t\in\mathbb{R}^n$), 因而它在每个连通分支上传递 (给定 $\alpha\in L_f$, 可写成 $\beta=A(t,\alpha)$, $t\in\mathbb{R}^n$ 的点 $\beta\in L_f$ 所成的集合是开集, 不能如此表示的点所成的集合也是开集).
\end{RiemannProof}

上述作用的稳定子群是 $\mathbb{R}^n$ 的离散子群. 下面描述这种子群的结构.

\begin{RiemannProposition}\label{prop:discrete-subgroups-of-euclidean-space}
设 $\Gamma$ 为 $\mathbb{R}^n$ 的离散子群. 则存在 $k\in\{0,1,\ldots,n\}$ 以及线性无关的向量 $e_1,\ldots,e_k$, 使得 $\Gamma=\operatorname{span}_{\mathbb{Z}}\{e_1,\ldots,e_k\}$.
\end{RiemannProposition}

\begin{RiemannProof}
若 $\Gamma=\{0\}$, 结论即得. 否则取 $e\in\Gamma\setminus\{0\}$. 由于 $\Gamma$ 离散, 集合
\[
  \Gamma\cap\{\lambda e\mid 0<\lambda\leq1\}
\]
有限 (且非空). 令 $e_1$ 为此集合中最接近 $0$ 的元素. 则
\[
  \Gamma\cap\operatorname{span}_{\mathbb{R}}\{e\}
  =\operatorname{span}_{\mathbb{Z}}\{e_1\}
\]
(否则, $e_1$ 就不是此集合中最接近 $0$ 的元素). 如果 $\Gamma=\operatorname{span}_{\mathbb{Z}}\{e_1\}$, 结论即得. 否则取 $e\in\Gamma\setminus\operatorname{span}_{\mathbb{Z}}\{e_1\}$. 则集合
\[
  \Gamma\cap\{\lambda e+\lambda_1e_1\mid 0<\lambda,\lambda_1\leq1\}
\]
有限 (且非空). 令 $e_2$ 为此集合中最接近 $\operatorname{span}_{\mathbb{R}}\{e_1\}$ 的元素 (Fig.~\ref{fig:discrete-subgroup-generators}). 则
\[
  \Gamma\cap\operatorname{span}_{\mathbb{R}}\{e,e_1\}
  =\operatorname{span}_{\mathbb{Z}}\{e_1,e_2\}.
\]
重复这一过程即得结论.
\end{RiemannProof}
\RGfigure[.46\textwidth]{discrete-subgroup-generators}{Proposition~\ref{prop:discrete-subgroups-of-euclidean-space} 的证明}

\begin{RiemannProposition}\label{prop:topology-of-invariant-level-components}
设 $L_f^\alpha$ 是 $\alpha\in L_f$ 所在的连通分支. 则 $L_f^\alpha$ 与 $T^k\times\mathbb{R}^{n-k}$ 微分同胚, 其中 $k$ 是稳定子群 $\Gamma_\alpha$ 的生成元个数. 特别地, 若 $L_f^\alpha$ 紧, 则它与 $n$ 维环面 $T^n$ 微分同胚.
\end{RiemannProposition}

\begin{RiemannProof}
由于作用 $A:\mathbb{R}^n\times L_f^\alpha\to L_f^\alpha$ 是传递的, 局部微分同胚 $A_\alpha:\mathbb{R}^n\to L_f^\alpha$ 是满射. 另一方面, 由于 $\Gamma_\alpha$ 离散, $\Gamma_\alpha$ 在 $\mathbb{R}^n$ 上的平移作用是自由且适当的, 因而可取商 $\mathbb{R}^n/\Gamma_\alpha$. 该商显然与 $T^k\times\mathbb{R}^{n-k}$ 微分同胚. 最后, 容易看出 $A_\alpha$ 诱导一个微分同胚 $\mathbb{R}^n/\Gamma_\alpha\cong L_f^\alpha$.
\end{RiemannProof}

现在可以理解完全可积系统上的 Hamilton 流了. 为此需要下面的定义 (参见 Fig.~\ref{fig:linear-flow-on-two-torus}).

\begin{RiemannDefinition}\label{def:linear-flow-and-frequencies-on-torus}
环面 $T^n=\mathbb{R}^n/\mathbb{Z}^n$ 上的\textbf{线性流}是由
\[
  \psi_t(x)=x+\nu t
\]
给出的流 $\psi_t:\mathbb{R}^n\to\mathbb{R}^n$ 的投影. 该线性流的\textbf{频率}是 $\nu$ 的各分量 $\nu^1,\ldots,\nu^n$.
\end{RiemannDefinition}
\RGfigure[.39\textwidth]{linear-flow-on-two-torus}{2 维环面上的线性流}

\begin{RiemannTheorem}[Arnold–Liouville]\label{thm:arnold-liouville-for-integrable-hamiltonian-systems}
设 $H$ 为完全可积 Hamilton 函数, 具有在稠密开集 $U\subset T^*M$ 上独立的 $n$ 个对合第一积分 $F_1,\ldots,F_n\in C^\infty(T^*M)$. 若映射 $(F_1,\ldots,F_n):U\to\mathbb{R}^n$ 的水平集的各连通分支都是紧的, 则它们是 $n$ 维环面, 且在 $X_H$ 的流下不变. $X_H$ 在这些环面上的流是线性流 (在适当选取的坐标下).
\end{RiemannTheorem}

\begin{RiemannProof}
由于 $(F_1,\ldots,F_n)$ 的水平集的各连通分支都是紧的, 向量场 $X_{F_i}$ 完备. 只需再证明 $X_H$ 在不变环面上的流是线性流. 在 Proposition~\ref{prop:topology-of-invariant-level-components} 给出的坐标下, 每个 $X_{F_i}$ 的流显然都是线性的. 由于 $X_H$ 与不变环面相切, 对某些函数 $f^i$ 有 $X_H=\sum_{i=1}^n f^iX_{F_i}$. 现在
\[
  0=X_{\{F_i,H\}}=[X_{F_i},X_H]=\sum_{j=1}^n(X_{F_i}\cdot f^j)X_{F_j},
\]
所以每个函数 $f^i$ 都在不变环面上恒定. 由此可知 $X_H$ 的流是线性的.
\end{RiemannProof}

接下来详细探究环面上线性流的性质.

\begin{RiemannDefinition}\label{def:time-average-along-linear-torus-flow}
设 $\psi_t:T^n\to T^n$ 为线性流. 函数 $f\in C^\infty(T^n)$ 沿 $\psi_t$ 的\textbf{时间平均}是映射
\[
  \overline f(x):=\lim_{T\to+\infty}\frac{1}{T}\int_0^T f(\psi_t(x))\,\mathrm{d}t
\]
(定义在使该极限存在的点 $x\in T^n$ 所成的集合上).
\end{RiemannDefinition}

\begin{RiemannDefinition}\label{def:independent-frequencies-for-linear-flow}
若线性流 $\psi_t:T^n\to T^n$ 的频率 $\nu\in\mathbb{R}^n$ 在 $\mathbb{Q}$ 上线性无关, 即对所有 $k\in\mathbb{Z}^n\setminus\{0\}$ 都有 $\langle k,\nu\rangle\ne0$, 则称其频率\textbf{独立}.
\end{RiemannDefinition}

\begin{RiemannTheorem}[Birkhoff 遍历性]\label{thm:birkhoff-ergodicity-for-linear-torus-flow}
若线性流 $\psi_t:T^n\to T^n$ 的频率 $\nu\in\mathbb{R}^n$ 独立, 则任意函数 $f\in C^\infty(T^n)$ 的时间平均对所有 $x\in T^n$ 都存在, 且
\[
  \overline f(x)=\int_{T^n}f.
\]
\end{RiemannTheorem}

\begin{RiemannProof}
由于 $T^n=\mathbb{R}^n/\mathbb{Z}^n$, 环面上的可微函数来自 $\mathbb{R}^n$ 上的周期可微函数, 而后者可以展开为一致收敛的 Fourier 级数. 因此只需证明定理对 $f(x)=e^{2\pi i\langle k,x\rangle}$ ($k\in\mathbb{Z}^n$) 成立.

若 $k=0$, 则等式两边都等于 $1$, 定理成立.

若 $k\ne0$, 则等式右端为零, 而左端为
\begin{align*}
  \overline f(x)
  &=\lim_{T\to+\infty}\frac{1}{T}\int_0^T
      e^{2\pi i\langle k,x+\nu t\rangle}\,\mathrm{d}t\\
  &=\lim_{T\to+\infty}\frac{1}{T}e^{2\pi i\langle k,x\rangle}
      \frac{e^{2\pi i\langle k,\nu\rangle T}-1}{2\pi i\langle k,\nu\rangle}=0
\end{align*}
(这里用到了 $\langle k,\nu\rangle\ne0$).
\end{RiemannProof}

\begin{RiemannCorollary}\label{cor:forward-linear-torus-orbits-dense}
若线性流 $\psi_t:T^n\to T^n$ 的频率独立, 则对所有 $x\in T^n$, $\{\psi_t(x)\mid t\geq0\}$ 在环面上稠密.
\end{RiemannCorollary}

\begin{RiemannProof}
若 $\{\psi_t(x)\mid t\geq0\}$ 不稠密, 则存在一个与它不相交的开集 $U\subset T^n$. 因而, 对任意支集非空且包含在 $U$ 中的非负函数 $f\in C^\infty(T^n)$, 都有 $\overline f(x)=0$ 而 $\int_{T^n}f>0$, 这与 Birkhoff 遍历性定理矛盾.
\end{RiemannProof}

\begin{RiemannCorollary}\label{cor:aperiodic-linear-torus-flow}
若线性流 $\psi_t:T^n\to T^n$ 的频率独立, 且 $n\geq2$, 则 $\psi_t(x)$ 不是周期的.
\end{RiemannCorollary}

\begin{RiemannRemark}\label{rem:kam-theorem-and-resonant-tori}
完全可积 Hamilton 函数生成的 Hamilton 流的定性行为已被完全理解. 然而完全可积性是一个非常强的条件, 一般的 Hamilton 函数并不满足该条件. \textbf{Kolmogorov–Arnold–Moser (KAM) 定理}给出了某种有限的普遍性: 当 Hamilton 函数受到小扰动时, 完全可积 Hamilton 函数的不变环面中有很大一部分仍然存在, 这些环面上的流仍是线性的, 且频率不变. 另一方面, 许多不变环面通常会被破坏, 其中包括那些频率不独立的环面 (\textbf{共振环面}).
\end{RiemannRemark}

\begin{RiemannExercise}\label{ex:completely-integrable-systems}
\begin{enumerate}[label=(\arabic*),ref=\arabic*,leftmargin=*]
\item\label{ex:poisson-bracket-of-first-integrals}
证明若 $F,G\in C^\infty(T^*M)$ 为第一积分, 则 $\{F,G\}$ 也是第一积分.
\item\label{ex:prove-involutive-integral-dimension-bound}
证明 Proposition~\ref{prop:dimension-bound-for-independent-integrals-in-involution}.
\item\label{ex:complete-integrability-of-surface-of-revolution}
考虑旋转曲面 $M\subset\mathbb{R}^3$, 它在柱坐标 $(r,\theta,z)$ 下由
\[
  r=f(z)
\]
给出, 其中 $f:(a,b)\to(0,+\infty)$ 可微.
\begin{enumerate}[label=(\alph*),leftmargin=*]
\item 证明 $M$ 的测地线是 Lagrange 函数 $L:TM\to\mathbb{R}$ 所确定的作用量的临界点, 其中 $L$ 在局部坐标下为
\[
  L(\theta,z,v^\theta,v^z)
    =\frac{1}{2}\bigl((f(z))^2(v^\theta)^2+((f'(z))^2+1)(v^z)^2\bigr).
\]
\item 证明局部坐标下由 $\theta=\text{常数}$ 或 $f'(z)=0$ 给出的曲线是测地线的像.
\item 求 Legendre 变换, 证明 $L$ 超正则, 并写出 Hamilton 函数 $H:T^*M\to\mathbb{R}$ 在局部坐标下的表达式.
\item 证明 $H$ 完全可积.
\item 证明不变集
\[
  L_{(E,l)}:=H^{-1}(E)\cap p_\theta^{-1}(l)
\]
($E,l>0$) 在 $M$ 上的投影在局部坐标下由
\[
  f(z)\geq\frac{l}{\sqrt{2E}}
\]
给出. 利用这一事实推出: 若 $f$ 在 $z=z_0$ 处有严格局部极大值, 则像为 $z=z_0$ 的测地线\textbf{稳定}, 即初始条件接近 $TM$ 中坐标为 $(\theta_0,z_0,1,0)$ 的点的测地线, 始终靠近曲线 $z=z_0$.
\end{enumerate}
\item\label{ex:stability-of-central-field-circular-orbits}
回顾 Example~\ref{ex:complete-integrability-in-dimension-one-and-central-fields}: 质量为 $m>0$ 的粒子在中心力场中的运动, 由完全可积的 Hamilton 函数
\[
  H(r,\theta,p_r,p_\theta)=\frac{p_r^2}{2m}+\frac{p_\theta^2}{2mr^2}+u(r)
\]
描述.
\begin{enumerate}[label=(\alph*),leftmargin=*]
\item 证明只要 $u'(r_0)\geq0$, 就存在半径为 $r_0$ 的圆轨道.
\item 核实 $\mathrm{d}H$ 与 $\mathrm{d}p_\theta$ 不独立的点所成的集合, 就是这些圆轨道的并.
\item 证明不变集
\[
  L_{(E,l)}:=H^{-1}(E)\cap p_\theta^{-1}(l)
\]
在 $\mathbb{R}^2$ 上的投影在局部坐标下由
\[
  u(r)+\frac{l^2}{2mr^2}\leq E
\]
给出.
\item 推出若 $u'(r_0)\geq0$ 且
\[
  u''(r_0)+\frac{3u'(r_0)}{r_0}>0,
\]
则半径为 $r_0$ 的圆轨道稳定.
\end{enumerate}
\item\label{ex:relativistic-circular-orbits}
在广义相对论中, 粒子在质量为 $M>0$ 的质点的引力场中的运动, 由在柱坐标 $(u,r,\theta)$ 下写出的 Lagrange 函数 $L:TU\to\mathbb{R}$ 给出:
\[
  L=-\frac12\left(1-\frac{2M}{r}\right)(v^u)^2
    +\frac12\left(1-\frac{2M}{r}\right)^{-1}(v^r)^2
    +\frac12 r^2(v^\theta)^2,
\]
其中 $U\subset\mathbb{R}^3$ 是由 $r>2M$ 给出的开集 (坐标 $u$ 称为\textbf{时间坐标}, 一般不同于粒子的\textbf{固有时}, 即曲线的参数 $t$).
\begin{enumerate}[label=(\alph*),leftmargin=*]
\item 证明 $L$ 超正则, 并求相应的 Hamilton 函数 $H:T^*U\to\mathbb{R}$.
\item 证明 $H$ 完全可积.
\item 证明对任意半径 $r_0>2M$, 都存在相应的圆轨道, 且当 $r_0>3M$ 时 $H<0$, 当 $r_0=3M$ 时 $H=0$, 当 $r_0<3M$ 时 $H>0$. (\textbf{注}: $H>0$ 的轨道不符合物理, 因为它们对应的速率超过光速; $H=0$ 的轨道只可能由以光速运动的无质量粒子实现.)
\item 证明满足 $\mathrm{d}H$, $\mathrm{d}p_u$ 和 $\mathrm{d}p_\theta$ 不独立 (且 $p_u\ne0$) 的点所成的集合, 就是这些圆轨道的并.
\item 证明不变柱面
\[
  L_{(E,k,l)}:=H^{-1}(E)\cap p_u^{-1}(k)\cap p_\theta^{-1}(l)
\]
在 $U$ 上的投影在局部坐标下由
\[
  \frac{l^2}{r^2}-\left(1-\frac{2M}{r}\right)^{-1}k^2\leq2E
\]
给出.
\item 推出若 $r_0>6M$, 则半径为 $r_0$ 的圆轨道稳定.
\end{enumerate}
\item\label{ex:complete-integrability-of-lagrange-top}
回顾, Lagrange 陀螺是由 Lagrange 函数 $L:TSO(3)\to\mathbb{R}$ 确定的力学系统, 其中 $L$ 在局部坐标下为
\[
  L=\frac{I_1}{2}\bigl((v^\theta)^2+(v^\varphi)^2\sin^2\theta\bigr)
     +\frac{I_3}{2}(v^\psi+v^\varphi\cos\theta)^2-Mgl\cos\theta,
\]
其中 $(\theta,\varphi,\psi)$ 为 Euler 角, $M$ 为陀螺的质量, $l$ 为定点到质心的距离.
\begin{enumerate}[label=(\alph*),leftmargin=*]
\item 求 Legendre 变换, 证明 $L$ 超正则, 并写出 Hamilton 函数 $H:T^*SO(3)\to\mathbb{R}$ 在局部坐标下的表达式.
\item 证明 $H$ 完全可积.
\item 证明 Exercise~\ref{ex:rigid-body}(\ref{ex:lagrange-top}) 中找到的满足 $|\dot\varphi|\ll|\dot\psi|$ 的解, 当 $|\dot\psi|$ 足够大时是稳定的.
\end{enumerate}
\item\label{ex:complete-integrability-of-euler-top}
证明满足 $I_1<I_2<I_3$ 的 Euler 陀螺在 $T^*SO(3)$ 上确定一个完全可积的 Hamilton 函数.
\item\label{ex:leading-digits-of-powers-of-two}
考虑由整数 $2^n$ ($n\in\mathbb{N}_0$) 的十进制展开的首位数字组成的数列:
\[
  1,2,4,8,1,3,6,1,2,5,1,2,4,8,1,3,6,1,2,5,\ldots
\]
本题的目的是回答: 这个数列中有 $7$ 吗?
\begin{enumerate}[label=(\alph*),leftmargin=*]
\item 证明若 $\nu\in\mathbb{R}\setminus\mathbb{Q}$, 则
\[
  \lim_{n\to+\infty}\frac{1}{n+1}\sum_{k=0}^n e^{2\pi i\nu k}=0.
\]
\item \phantomsection\label{idx:discrete-birkhoff}证明 Birkhoff 遍历性定理的如下离散版本: 若可微函数 $f:\mathbb{R}\to\mathbb{R}$ 的周期为 $1$, 且 $\nu\in\mathbb{R}\setminus\mathbb{Q}$, 则对所有 $x\in\mathbb{R}$ 有
\[
  \lim_{n\to+\infty}\frac{1}{n+1}\sum_{k=0}^n f(x+\nu k)=\int_0^1 f(x)\,\mathrm{d}x.
\]
\item 证明 $\log 2$ 是 $\log 10$ 的无理数倍.
\item 上面的数列中有 $7$ 吗?
\end{enumerate}
\end{enumerate}
\end{RiemannExercise}
